% Appendix A: formal foundations, revised August 2026
\label{app:mathematical_foundations}

This appendix states exactly what the H-Potential mathematics does and does not
establish. It locates the construction in preceding mathematical traditions,
defines its objects and scales, states axioms before selecting a functional
form, and separates mathematical consequences from empirical hypotheses. Every
formal statement is followed by a plain-language interpretation.

\begin{ataglance}[How to read this appendix]
The appendix is an argument, not a reference list, and it runs in one
direction. It first borrows: measurement theory, utility representation,
production functions and multidimensional deprivation each contribute a
technique, and none contributes a result about health
(\S\ref{app:a-preceding}). It then fixes what is being measured and when
(\S\ref{app:a-primitives}). Only then does it state the properties a summary
score must have (\S\ref{app:a-axioms}), and show that the most common choice,
adding the domains up, quietly violates one of them
(\S\ref{app:a-additivity}). The CES family that survives the axioms is
developed next (\S\ref{app:a-ces}), followed by a rival that denies one of the
axioms outright (\S\ref{app:a-relational}) and a threshold alternative that
denies smoothness (\S\ref{app:a-threshold}). The last sections separate the
construct from its indicators, list the objects that get confused, state the
conditions under which the whole framework fails, and connect the score to
money.

A reader who wants the argument without the proofs can read the
\emph{In plain language} remark after each result; they are written to stand
alone. The two reference profiles of
\S\ref{sec:intro-cases}, Profile~A and Profile~B, are carried through
\S\ref{app:a-axioms}--\S\ref{app:a-threshold} with numbers attached.
\end{ataglance}

\section*{The Appendix in Plain Terms}

The main chapters build a score from three domains and test whether it
predicts death and hospital admission. This appendix does the groundwork that
a laboratory does before it reports a new assay. It states what is measured, on
what scale, what the reported number may and may not be taken to mean, and
which findings would show the method unfit for use. Its results are proofs:
statements that follow from stated assumptions by logic alone. They use no
NHATS data, and none of the findings of Chapters~\ref{ch:results1}
and~\ref{ch:results2} depends on the reader checking them line by line.

Three points carry into the main argument. First, the score orders people but
has no natural unit. An assay reported in arbitrary units is in the same
position until it is calibrated against an international standard, and the
factor score of Chapter~\ref{ch:methods} (Definition~\ref{def:factorscore}) is
reported on that understanding. Second, the rule that combines the three
domains is a choice with ethical content. Adding the domains lets a gain in one
cancel a loss in another; the multiplicative rules of Chapter~\ref{app:crf}
limit that trade; and two reasonable rules can put the same two people in
opposite order. For this reason Chapter~\ref{ch:methods} (\S\ref{sec:roadmap})
requires every scalar result to be reported with its sensitivity to the rule.
Third, the person's profile, the chosen score, the estimated score and the
predicted risk are four different quantities, and none of them settles whether
an intervention is good value, or whether a funding decision made in one health
system holds in another. Chapter~\ref{ch:theory} (\S\ref{sec:lit-economics})
and Chapter~\ref{app:crf} (Theorem~\ref{thm:nontransport}) rest on these last
results.

Table~\ref{tab:A-results} lists each numbered result, states it in one line
and shows where the main text uses it. A reader who accepts the one-line
statements can skip the proofs. Three terms recur. A \emph{theorem} or
\emph{proposition} is a statement proved from the definitions. An
\emph{axiom} is a requirement adopted by choice and not proved; it plays the
part of an acceptance criterion fixed before a method validation. A
\emph{counterexample} is a single case that shows a general claim to be false,
and one is enough.

\begin{table}[htbp]\centering
\footnotesize
\caption{The numbered results of Appendix~\ref{app:mathematical_foundations}, what each says in plain words, and where the main text uses it.}
\label{tab:A-results}
\begin{tabular}{p{3.4cm}p{5.8cm}p{3.9cm}}
\toprule
Result (section) & Plain statement & Used in \\
\midrule
Dominance preservation (\S\ref{app:a-axioms}) & A profile at least as good in every domain is never scored lower. & Ch.~\ref{app:crf}, \S\ref{sec:E-axioms} \\
Ordinal representation is not cardinal measurement (\S\ref{app:a-axioms}) & A score that orders people does not say how far apart they are; gaps and ratios need separate calibration. & Ch.~\ref{ch:methods}, \S\ref{sec:roadmap}; Ch.~\ref{ch:theory}, \S\ref{sec:lit-economics} \\
Compensation (\S\ref{app:a-additivity}) & Adding the domains lets a gain in one cancel a loss in another, at one fixed rate. & Ch.~\ref{ch:methods}, \S\ref{sec:roadmap}; Ch.~\ref{ch:theory}, \S\ref{sec:lit-candidates} \\
Admissibility (\S\ref{app:a-ces}) & The CES score lies between zero and biomedical capacity and never falls when a domain improves. & Ch.~\ref{ch:theory}, Prop.~\ref{prop:lit-heuristic}; Ch.~\ref{app:crf}, \S\ref{sec:E-family} \\
Constant substitution elasticity (\S\ref{app:a-ces}) & One parameter, $\rho$, fixes how readily social enablement and structural freedom stand in for each other. & Ch.~\ref{ch:methods}, \S\ref{sec:roadmap}; Ch.~\ref{app:crf}, \S\ref{sec:E-family} \\
Geometric marginal effects (\S\ref{app:a-ces}) & In the geometric case each weight is the percentage change in the score for a one per cent change in that domain. & Ch.~\ref{app:crf}, \S\ref{sec:E-marginal} \\
Capacity--context complementarity (\S\ref{app:a-ces}) & Treating the body yields more where the context is better, and the reverse: an interaction, not two separate effects. & Ch.~\ref{ch:theory}, Assumption~\ref{ass:lit-interaction}; Ch.~\ref{app:crf}, \S\ref{sec:E-marginal} \\
Aggregation can reverse rankings (\S\ref{app:a-ces}) & Two admissible rules can put the same two people in opposite order. & Ch.~\ref{ch:theory}, Prop.~\ref{prop:lit-reversal}; Ch.~\ref{ch:methods}, \S\ref{sec:roadmap} \\
Properties of the relational model (\S\ref{app:a-relational}) & The rival score is well behaved but can be above zero when biomedical capacity is zero. & Ch.~\ref{ch:methods}, \S\ref{sec:roadmap}; Ch.~\ref{app:crf}, \S\ref{sec:E-ontologies} \\
The ontology is not selected by fit alone (\S\ref{app:a-relational}) & Better fit to data does not show which of the two views of health is right. & Ch.~\ref{ch:methods}, \S\ref{sec:roadmap}; Ch.~\ref{app:crf}, \S\ref{sec:E-ontologies} \\
Latent scale indeterminacy, Thm.~\ref{thm:latent_scale} & The zero and unit of an estimated domain score are set by convention. & Ch.~\ref{ch:methods}, Def.~\ref{def:factorscore} \\
Second-order implied covariance, Prop.~\ref{prop:secondorder_cov} & The pattern of correlations among indicators that one higher-order factor predicts. & Ch.~\ref{ch:methods}, Def.~\ref{def:secondorder}; Ch.~\ref{ch:results1}, \S\ref{sec:anal-secondorder} \\
Non-equivalence (\S\ref{app:a-confusion}) & Profile, chosen score, estimated score and predicted risk cannot be swapped for one another. & Ch.~\ref{ch:methods}, \S\ref{sec:roadmap} \\
Need does not imply cost-effectiveness (\S\ref{app:a-economics}) & A low score shows need; it does not show that any intervention is good value. & Ch.~\ref{ch:theory}, \S\ref{sec:lit-economics} \\
A policy decision is not generally transportable (\S\ref{app:a-economics}) & The same effect in two health systems can justify funding in one and not the other. & Ch.~\ref{ch:theory}, \S\ref{sec:lit-economics}; Ch.~\ref{app:crf}, Thm.~\ref{thm:nontransport}; Ch.~\ref{ch:discussion}, \S\ref{sec:disc-impl-policy} \\
\bottomrule
\end{tabular}
\end{table}

\section{Preceding Theorems and Mathematical Spaces}
\label{app:a-preceding}

Measurement theory distinguishes nominal, ordinal, interval and ratio scales
and restricts the transformations under which a statement remains meaningful.
\citep{stevens_1946} Utility-representation theory supplies conditions under
which an ordering can be represented by a real-valued function. Debreu's
representation result requires completeness, transitivity, continuity and
suitable topological conditions; it does not establish cardinal differences,
interpersonal comparability or an additive form. \citep{debreu_1954} Additive
multi-attribute representations require stronger preferential-independence
conditions. \citep{keeney_raiffa_1976}

Production-function theory studies how inputs combine. The CES family was
introduced to represent different degrees of substitutability.
\citep{arrow_chenery_1961} Grossman's health-capital model treats health as a
durable stock affecting healthy time and influenced by investment.
\citep{grossman_1972} These models motivate questions about complementarity and
substitution; they do not prove that social support and biological capacity are
literal production inputs.

Axiomatic multidimensional measurement demonstrates that aggregation encodes
ethical choices. Alkire and Foster use dimension-specific thresholds and a
cross-dimensional cut-off, providing a transparent non-compensatory alternative
to a smooth total. \citep{alkire_foster_2011} Latent-variable theory addresses a
different problem: whether covariance among indicators can be explained by
unobserved variables. It requires an explicit direction between construct and
indicators. \citep{bollen_lennox_1991,jarvis_2003}

\begin{remark}[What is inherited]
The thesis inherits questions and proof techniques from these traditions. It
does not inherit a theorem stating that H-Potential exists as a unique cardinal
health quantity. That would require assumptions and evidence not supplied by
the preceding literature.
\end{remark}

What the borrowed traditions do supply is discipline about order. Each of them
begins by saying precisely what is being measured, on what scale, and at what
moment, before anything is combined. That is where this appendix begins too.

\section{Population, Time and Primitive Quantities}
\label{app:a-primitives}

The first definition fixes who is measured and when. As in any cohort study,
the baseline is recorded before any outcome is counted.

\begin{definition}[Population and index time]
Let $(\Omega,\mathcal F,\mathbb P)$ be a probability space. A person is
represented by $\omega\in\Omega$ and baseline occurs at $t_0$. Construct
indicators are measured at or before $t_0$; prognostic outcomes occur after
$t_0$.
\end{definition}

\begin{remark}[In plain language]
People differ, and a participant is treated as one possible draw from the
target population. The time rule prevents future death or hospitalisation from
being used to define the baseline construct.
\end{remark}

The next definition sets out the three numbers recorded for each person, each
rescaled to run from $0$ (least favourable) to $1$ (most favourable), much as
factor activity is reported as a percentage of pooled normal plasma rather than
as a raw clotting time.

\begin{definition}[Domain profile]
The normalized profile is
\[
X(\omega)=(P(\omega),S(\omega),R(\omega))\in\mathcal X=[0,1]^3,
\]
where $P$ is biomedical and functional capacity, $S$ social and environmental
enablement, and $R=1-B$ freedom from structural burden. Larger values are
favourable. Normalization is part of the protocol, not a claim that raw
variables possess ratio scales.
\end{definition}

\begin{remark}[In plain language]
The profile $(0.8,0.4,0.9)$ records relatively high capacity, limited
enablement and relatively little burden. It does not yet combine these facts,
and $0.8$ is not asserted to be literally twice $0.4$.
\end{remark}

\begin{workedexample}[apa-profiles]{The two reference profiles}
Profile~A (Table~\ref{tab:profiles}) has moderate capacity, low enablement and
low structural freedom: $X=(0.55,0.35,0.40)$. Profile~B has low capacity and
high enablement. With structure favourable it is $(0.30,0.70,0.85)$; with
structure unfavourable it is $(0.30,0.70,0.20)$, with only the third coordinate
changed.

Under the geometric rule of \S\ref{app:a-ces} with $\alpha=0.5$, these give
$0.55\sqrt{0.35\times0.40}=0.206$ for Profile~A, $0.30\sqrt{0.70\times0.85}=0.231$
for Profile~B with structure favourable and $0.30\sqrt{0.70\times0.20}=0.112$
with structure unfavourable. The structural coordinate alone accounts for a
factor of two in the score.

\begin{wecode}[Python]
from hpv.crf import H

cases = {"Profile A": (0.55, 0.35, 0.40),
         "Profile B favourable": (0.30, 0.70, 0.85),
         "Profile B unfavourable": (0.30, 0.70, 0.20)}
for name, x in cases.items():
    print(f"{name:<22} profile {x}  H = {H(*x):.3f}")
print("favourable / unfavourable =", H(0.30, 0.70, 0.85) / H(0.30, 0.70, 0.20))
\end{wecode}

\weanswer{$0.206$, $0.231$ and $0.112$; the ratio of favourable to
unfavourable is $2.06$. Asserted by
\texttt{test\_running\_cases\_reproduce\_the\_thesis\_numbers}.}
\end{workedexample}

\begin{reading}[What a profile does not settle]
A profile does not say whether Profile~A is worse off than Profile~B, because
$(0.55,0.35,0.40)$ and $(0.30,0.70,0.85)$ are not comparable by dominance: each is better in some coordinate. Producing an answer
requires an extra assumption, and the rest of the appendix is about which extra
assumptions are defensible.
\end{reading}

One comparison between profiles needs no assumption at all, and the next
definition states it.

\begin{definition}[Dominance]
For $x,y\in\mathcal X$, write $x\succeq_Dy$ if $x_j\ge y_j$ in every
coordinate, and $x\succ_Dy$ if at least one inequality is strict.
\end{definition}

\begin{remark}[In plain language]
Dominance is the weakest comparison: one profile is no worse anywhere and
better somewhere. Profiles that trade strength in one domain against weakness
in another cannot be ranked by dominance alone.
\end{remark}

Dominance settles some comparisons and not others. Profile~B with structure
favourable dominates Profile~B with structure unfavourable, so the higher
structural freedom is better without further assumption. Profile~A and
Profile~B are not ranked by dominance. Since clinical and budgetary decisions
are made mostly among such profiles, something stronger than dominance is
required, and the next section states the least that it requires.

\section{Axioms for Any Scalar Score}
\label{app:a-axioms}

This section lists the least that any single summary score must satisfy. The
requirements are chosen in advance, as a laboratory fixes acceptance criteria
before it validates a method.

\begin{definition}[Admissible aggregator]
An aggregator $H:\mathcal X\to[0,1]$ is admissible for the primary sensitivity
family if it satisfies:
\begin{enumerate}
 \item continuity;
 \item weak monotonicity: $x\succeq_Dy\Rightarrow H(x)\ge H(y)$;
 \item normalization: $H(1,1,1)=1$ and $H(0,s,r)=0$;
 \item capacity boundedness: $0\le H(p,s,r)\le p$;
 \item transparency of weights and substitution assumptions.
\end{enumerate}
\end{definition}

\begin{remark}[In plain language]
These are design commitments. The boundary $H(0,s,r)=0$ says that context
cannot manufacture biomedical capacity from nothing. Readers who reject this
claim should reject this family or use a less restrictive alternative.
\end{remark}

\begin{reading}[What each axiom rules out, clinically]
\textbf{Continuity} says a hair's-breadth change in circumstance cannot flip
the score. It rules out a measure in which earning one rand above a means-test
line halves a person's recorded capability.

\textbf{Weak monotonicity} says that making something better cannot make the
score worse. It rules out an index in which an improvement to a domain lowers the score.

\textbf{Normalization} fixes the two ends of the ruler.

\textbf{Capacity boundedness}, $H\le p$, is the contentious one. It says that
Profile~B, with $P=0.30$, cannot be scored above the ceiling its biomedical
coordinate sets, however high $S$ and $R$ are. A favourable context can bring
the score close to that ceiling; it cannot lift it past $p$. Readers who hold
that an accessible environment is part of health, not a support for it, will
reject this axiom, and \S\ref{app:a-relational} is written for them.
\end{reading}

The first consequence follows at once from the monotonicity requirement.

\begin{proposition}[Dominance preservation]
Every weakly monotone aggregator preserves dominance.
\end{proposition}
The proof applies that requirement directly.
\begin{proof}
If $x\succeq_Dy$, weak monotonicity gives $H(x)\ge H(y)$.
\end{proof}

\begin{remark}[In plain language]
Improving at least one favourable domain while worsening none cannot lower the
score. This does not determine how to rank trade-offs.
\end{remark}

The next result limits what any such score can mean. A score that only orders
people carries no information about the size of the gaps between them.

\begin{theorem}[Ordinal representation is not cardinal measurement]
Suppose $H$ represents an ordering, so
$x\succeq y\Leftrightarrow H(x)\ge H(y)$. For every strictly increasing
$g$, $g\circ H$ represents the same ordering. Therefore differences and ratios
of scores are not identified by the ordering alone.
\end{theorem}
The proof stretches the scale in a way that keeps the order, for example by
taking a square root, and shows that the gaps change.
\begin{proof}
Strict monotonicity of $g$ gives
$H(x)\ge H(y)\Leftrightarrow g(H(x))\ge g(H(y))$. Yet nonlinear $g$
generally changes differences and ratios.
\end{proof}

\begin{remark}[In plain language]
A score of 80 is not proved to represent twice the potential of 40. A
ten-point change is not automatically equal everywhere on the scale. Years,
utilities or clinical units require an independent calibration study.
\end{remark}

Clinicians meet this constantly and mostly handle it correctly. Nobody claims a
patient with a NYHA class~IV heart failure is twice as breathless as class~II,
or that the gap between Glasgow Coma Scale 6 and 8 equals the gap between 12
and 14. The scale orders; it does not measure distance. The error the theorem
guards against is the one made downstream, when an ordinal score is fed into a
cost-effectiveness calculation that multiplies and divides it as though it were
a quantity of something.

If the ordering is all that is identified, then a great deal rides on how the
domains are combined into that ordering, because different combination rules
give different orderings. The obvious rule is to add them up, and it is worth
seeing exactly what that commits one to.

\section{Why Additivity Is a Strong Assumption}
\label{app:a-additivity}

The commonest way to combine domains is a weighted sum. The definition states
it so that its consequence can be proved.

\begin{definition}[Additive form]
\[
H_A(p,s,r)=w_Pu_P(p)+w_Su_S(s)+w_Ru_R(r),\qquad
w_j\ge0,\quad\sum_jw_j=1.
\]
\end{definition}

The next result shows what the weighted sum commits its user to.

\begin{proposition}[Compensation]
If $u_j(x)=x$ and $w_P,w_S>0$, a loss $\Delta p$ can be offset exactly by
$\Delta s=(w_P/w_S)\Delta p$, whenever the changed profile remains in
$\mathcal X$.
\end{proposition}
The proof substitutes the stated gain into the score.
\begin{proof}
The score change is $-w_P\Delta p+w_S\Delta s=0$ after substitution.
\end{proof}

\begin{remark}[In plain language]
An additive score allows social enablement to cancel lost capacity exactly.
That may or may not be ethically and clinically acceptable. Additivity is not
a neutral default, and its classical representation requires independence
conditions not guaranteed by fitting a regression. \citep{keeney_raiffa_1976}
\end{remark}

\begin{workedexample}[apa-additive]{An additive index can rise when structural freedom collapses}
Take equal weights, $w_P=w_S=w_R=1/3$, and Profile~B with structure
unfavourable: $(0.30,0.70,0.20)$. Its additive score is
$(0.30+0.70+0.20)/3=0.400$.

Now raise biomedical potential from $0.30$ to $0.60$ and, at the same time,
lower structural freedom from $0.20$ to $0.04$. The additive score becomes
$(0.60+0.70+0.04)/3=0.447$. It has gone \emph{up}.

The realisation function of \S\ref{app:a-ces} moves the other way:
$0.60\sqrt{0.70\times0.04}=0.100$ against $0.30\sqrt{0.70\times0.20}=0.112$, so
realised capability falls.

\begin{wecode}[Python]
from hpv.crf import H, H_additive

w = (1/3, 1/3, 1/3)
before, after = (0.30, 0.70, 0.20), (0.60, 0.70, 0.04)
for label, x in (("before", before), ("after", after)):
    print(f"{label:<22} additive {H_additive(*x, w):.4f}   CRF {H(*x):.4f}")
print("additive change", H_additive(*after, w) - H_additive(*before, w))
print("CRF change     ", H(*after) - H(*before))
\end{wecode}

\weanswer{The additive index rises from $0.400$ to $0.447$; the CRF falls from
$0.112$ to $0.100$. The two rules disagree about the sign of the change, which
is the sharpest form the disagreement can take. Asserted by
\texttt{test\_worked\_example\_additive\_perverse}.}
\end{workedexample}

\begin{reading}[Why the additive rule does this]
On the additive measure, a gain in biomedical potential paired with a collapse
in structural freedom counts as an improvement. The proposition above explains
why: with equal weights, one unit
of lost structural freedom is exactly cancelled by one unit of gained biology.
Additivity permits that trade and fixes its rate
everywhere on the scale, including at the bottom where the trade is least
defensible. Any measure intended to detect the effect of health systems on
people cannot use a rule that lets biology buy the system's failures back.
\end{reading}

The alternative is a form in which the domains multiply rather than add, so
that a near-zero in one cannot be bought off by an excess in another.

\section{CES Capability Realisation}
\label{app:a-ces}

This definition introduces the family used as the primary sensitivity
analysis. Biomedical capacity multiplies an average of the two contextual
domains, and one parameter sets what kind of average is taken.

\begin{definition}[CES contextual mean and realised potential]
For $0<\alpha<1$ and $\rho\le1$, define
\[
M_{\rho,\alpha}(s,r)=
\begin{cases}
[\alpha s^\rho+(1-\alpha)r^\rho]^{1/\rho},&\rho\ne0,\\
s^\alpha r^{1-\alpha},&\rho=0,
\end{cases}
\qquad
H_{\rho,\alpha}(p,s,r)=pM_{\rho,\alpha}(s,r),
\]
using the continuous extension at zero.
\end{definition}

\begin{remark}[In plain language]
Capacity sets the ceiling. The contextual mean describes how much can be
realised. The weight $\alpha$ controls relative emphasis; $\rho$ controls how
easily strength in one contextual domain compensates for weakness in another.
\end{remark}

\begin{reading}[Reading $\rho$ as a clinical question]
The parameter $\rho$ answers one question: how much can good social support
make up for a hostile system, and the reverse?

At $\rho=1$ the contextual mean is the ordinary weighted average, and the
answer is ``completely, at a fixed exchange rate''. At $\rho=0$ it is the
geometric mean, and the answer is ``a lot in the middle, almost nothing near
the extremes''. As $\rho\to-\infty$ the mean collapses to $\min(s,r)$ and the
answer is ``not at all: the weaker of the two is the whole story''.

The physiological analogue is oxygen delivery, which is cardiac output times
arterial oxygen content. A superb cardiac output cannot rescue a haemoglobin of
$2$, and a normal haemoglobin cannot rescue a cardiac arrest. The product
structure encodes the fact that both are necessary. Multidimensional
deprivation research reaches the same structure from ethics rather than
physiology, which is some reason to think the form is not merely convenient.
\end{reading}

\begin{workedexample}[apa-rho]{Three values of $\rho$, two profiles}
Set $\alpha=0.5$ and compute the contextual mean $M_{\rho,\alpha}(s,r)$ at the
three landmark values of $\rho$.

For Profile~A, whose contextual coordinates are balanced at $s=0.35$ and
$r=0.40$: the arithmetic mean ($\rho=1$) is $0.375$, the geometric mean
($\rho=0$) is $0.374$ and the bottleneck ($\rho\to-\infty$) is $0.350$. The
three answers sit within seven per cent of each other, so the choice of $\rho$
has little effect.

For Profile~B with structure unfavourable, whose contextual coordinates are
unbalanced at $s=0.70$ and $r=0.20$: arithmetic $0.450$, geometric $0.374$,
bottleneck $0.200$. Here the choice of $\rho$ changes the contextual mean by a
factor of $2.25$, and it
changes which of the two \emph{contexts} is ranked first.

\begin{wecode}[Python]
from hpv.crf import contextual_mean

for name, (s, r) in (("Profile A", (0.35, 0.40)), ("Profile B", (0.70, 0.20))):
    vals = [contextual_mean(s, r, rho) for rho in (1.0, 0.0, -800.0)]
    print(f"{name:<14} arithmetic {vals[0]:.3f}  geometric {vals[1]:.3f}  "
          f"bottleneck {vals[2]:.3f}  spread {max(vals)/min(vals):.2f}")
\end{wecode}

\weanswer{Profile~A $0.375$, $0.374$, $0.350$, a spread of $1.07$; Profile~B
with structure unfavourable $0.450$, $0.374$, $0.200$, a spread of $2.25$. On
the contextual mean the arithmetic rule ranks the context of Profile~B above
that of Profile~A, the bottleneck rule reverses the order, and the geometric
rule calls them equal. On realised capability $H$ Profile~A leads under all
three, because its biomedical potential is higher. Asserted by
\texttt{test\_worked\_example\_rho\_reverses\_the\_ranking} and
\texttt{test\_substitution\_ordering\_across\_rho}.}
\end{workedexample}

\noindent The parameter matters where the contextual coordinates are
unbalanced.

The next proposition checks that this family meets the requirements of
\S\ref{app:a-axioms}.

\begin{proposition}[Admissibility]
$H_{\rho,\alpha}$ is continuous, weakly increasing, normalized and satisfies
$0\le H_{\rho,\alpha}(p,s,r)\le p$.
\end{proposition}
The proof uses one fact: any such average lies between the smaller and the
larger of the two values averaged.
\begin{proof}
The weighted power mean is continuous, non-decreasing and bounded between
$\min(s,r)$ and $\max(s,r)$. Thus it lies in $[0,1]$. Multiplication by
$p\in[0,1]$ gives the bound and direct substitution gives normalization.
\end{proof}

\begin{remark}[In plain language]
The score never exceeds available capacity, and improving a favourable domain
cannot reduce it. These results follow from the chosen form. They do not prove
that CES is the true law of healthy ageing.
\end{remark}

The next result gives $\rho$ an exact meaning. The elasticity of substitution
measures how readily one input replaces the other.

\begin{proposition}[Constant substitution elasticity]
For interior values, the CES elasticity of substitution between $S$ and $R$ is
\citep{arrow_chenery_1961}
\[
\sigma=\frac{1}{1-\rho}.
\]
\end{proposition}
The proof is a standard calculation from production economics.
\begin{proof}
The ratio of marginal products is proportional to $(s/r)^{\rho-1}$.
Log-differentiating the input ratio with respect to this marginal-rate ratio
gives $1/(1-\rho)$.
\end{proof}

\begin{remark}[In plain language]
Values of $\rho$ near one allow substantial compensation; $\rho=0$ gives a
geometric mean \citep{hardy_littlewood_polya_1934}; very negative values make the weaker domain a bottleneck.
Ethically acceptable compensation cannot be selected by fit alone.
\end{remark}

For the geometric case the effect of a small change in each domain can be
written down exactly.

\begin{proposition}[Geometric marginal effects]
For $\rho=0$ and interior values,
\[
H=ps^\alpha r^{1-\alpha},\quad
H_p=H/p,\quad H_s=\alpha H/s,\quad H_r=(1-\alpha)H/r.
\]
The corresponding elasticities are $1,\alpha,1-\alpha$.
\end{proposition}
The proof is direct differentiation.
\begin{proof}
Differentiate the monomial and multiply each derivative by its argument divided
by $H$.
\end{proof}

\begin{remark}[In plain language]
The weights have a precise proportional interpretation in the geometric case.
They are modelling parameters, not natural constants.
\end{remark}

The next result states the central claim of the theory: the effect of a
change in capacity depends on the context, and the reverse. In the language of
a factorial trial, it is an interaction between two factors.

\begin{proposition}[Capacity--context complementarity]
For the geometric model,
\[
H_{ps}=\alpha s^{\alpha-1}r^{1-\alpha}>0,\qquad
H_{pr}=(1-\alpha)s^\alpha r^{-\alpha}>0.
\]
\end{proposition}
The proof differentiates the previous result once more.
\begin{proof}
Differentiate $H_p=s^\alpha r^{1-\alpha}$ with respect to $s$ and $r$.
\end{proof}

\begin{remark}[In plain language]
Additional capacity produces more realised benefit in a supportive context,
and contextual improvement produces more benefit when capacity exists. This is
the theory's central complementarity claim and must be tested empirically.
\end{remark}

\begin{reading}[Why the cross-partial is the intersectionality claim]
A positive $H_{ps}$ means the return on treating the body depends on the
circumstances the body is treated in. Put the other way round, the same
clinical intervention buys different amounts of realised capability in
different people, and the difference is systematic rather than random.

This is testable and it is the claim most likely to be wrong. It predicts, for
instance, that a given rise in $p$ yields less realised capability for
Profile~A than for a profile with the same $p$ and higher $S$ and $R$, and that
the shortfall is larger than either difference would predict alone. An additive
model predicts that the two effects add. The empirical
chapters test the interaction directly, and a null result would narrow the
theory and not simply fail to support it.
\end{reading}

The next result shows that the choice among admissible rules matters for
individual people, not just for averages.

\begin{proposition}[Aggregation can reverse rankings]
There are profiles ranked oppositely by additive and bottleneck rules.
\end{proposition}
The proof needs one pair of profiles on which two rules disagree.
\begin{proof}
Let $x=(1,1,0.2)$ and $y=(1,0.55,0.55)$. Equal-weight arithmetic contextual
means are $0.6$ and $0.55$, ranking $x$ first. Bottleneck values are $0.2$ and
$0.55$, ranking $y$ first.
\end{proof}

\begin{remark}[In plain language]
Whether excellence conceals a severe shortfall is not a technical detail.
Because rankings can reverse, weight and substitution sensitivity must be
reported.
\end{remark}

The counterexample in the proof is not exotic. Profile $x$ is a person with
strong support and a system that fails them; profile $y$ is a person with
moderate amounts of both. Which of the two a health service should prioritise
is a real question that gets a different answer depending on a parameter the
data cannot settle. Reporting the score without reporting its sensitivity to
$\rho$ would present an ethical choice as an empirical finding.

Everything so far has assumed the capacity ceiling. The next section takes that
assumption away and asks what changes.

\section{A Rival Relational Ontology}
\label{app:a-relational}

The CES ceiling model assumes $H(p,s,r)\le p$. Capability theory, assistive
technology and relational accounts of disability provide a reason to test a
different boundary.

The definition below gives a score without the capacity ceiling. Each domain,
and the product of capacity with each contextual domain, raises the score, and
the exponential form keeps it between $0$ and $1$.

\begin{definition}[Normalized relational aggregator]
For non-negative parameters
$\boldsymbol\beta=(\beta_P,\beta_S,\beta_R,\beta_{PS},\beta_{PR})$ with
$K=\beta_P+\beta_S+\beta_R+\beta_{PS}+\beta_{PR}>0$, define
\[
H_R(p,s,r)=
\frac{1-\exp[-(\beta_Pp+\beta_Ss+\beta_Rr+
\beta_{PS}ps+\beta_{PR}pr)]}{1-\exp(-K)}.
\]
\end{definition}

The next result shows that this rival meets every requirement except capacity
boundedness.

\begin{proposition}[Properties of the relational model]
$H_R$ is continuous, maps $[0,1]^3$ into $[0,1]$, is weakly increasing in every
coordinate and satisfies $H_R(1,1,1)=1$. If
$\beta_Ss+\beta_Rr>0$, then $H_R(0,s,r)>0$, so the capacity-ceiling axiom is
not satisfied.
\end{proposition}
The proof checks each property in turn.
\begin{proof}
The exponent argument is continuous, non-negative and no larger than $K$.
Therefore the numerator lies between zero and $1-e^{-K}$. Each partial
derivative is the positive exponential factor multiplied by a sum of
non-negative coefficients. Substitution establishes the two boundary claims.
\end{proof}

\begin{remark}[In plain language]
The relational model allows assistance, technology and institutions to create
effective opportunity not captured by an unaided biomedical-capacity measure.
This is philosophically closer to relational capability, but it also risks
calling general resources health. The two models embody rival views and must be
reported as such.
\end{remark}

\begin{reading}[Where the two ontologies disagree about Profile~B]
Under the ceiling model, the biomedical coordinate of Profile~B, $P=0.30$, sets
a limit that only a biomedical change can raise. Assistive technology,
accessible housing and supportive employment, recorded in $S$ and $R$, bring
realised capability closer to that limit; they cannot pass it.

Under the relational model they can. If assistive technology provides the
functionings that biomedical capacity would otherwise provide, then on this
account the health-related capability of Profile~B is higher, and not only
better supported. The disability-rights literature has argued this position for
forty years.

The two models make different predictions about a specific case: a person with
$p=0$ on every biomedical indicator. The ceiling model scores them zero by
construction. The relational model scores them by their circumstances. Neither
answer is obviously right, and no amount of NHATS data settles it, which is
what the following proposition states formally.
\end{reading}

\begin{proposition}[The ontology is not selected by fit alone]
If the ceiling and relational models are non-nested and encode different
boundary conditions, superior sample fit of one does not logically prove its
boundary condition to be philosophically correct.
\end{proposition}
The argument rests on the difference between describing data and fixing what
a word means.
\begin{proof}
Fit compares discrepancies between model-implied and observed distributions
under the available indicators. A boundary claim concerns the meaning of
capability under profiles that may be sparsely observed or imperfectly
measured. Descriptive adequacy on one sample does not entail the truth of that
normative or ontological premise.
\end{proof}

\begin{remark}[In plain language]
Data can show which model describes observed patterns better. They cannot, by
themselves, decide whether an accessible environment is part of health or a
separate condition supporting health.
\end{remark}

\section{A Non-Compensatory Threshold Alternative}
\label{app:a-threshold}

This section turns to a rule that never trades one domain against another. It
counts the domains that fall below a declared minimum, as a diagnostic rule
counts the criteria a patient meets.

\begin{definition}[Deprivation count]
Choose thresholds $z_j\in(0,1)$ and define
$d_j(x)=\mathbf1\{x_j<z_j\}$ and $c(x)=\sum_jw_jd_j(x)$. A profile is
classified as constrained when $c(x)\ge k$ for a declared cut-off $k$.
\end{definition}

\begin{remark}[In plain language]
This asks whether minimum standards are missed instead of allowing every
shortfall to be smoothed away. It adapts the logic of multidimensional poverty
measurement. \citep{alkire_foster_2011} Thresholds require clinical and public
justification.
\end{remark}

\begin{workedexample}[apa-deprivation]{Counting deprivations for Profile~A}
Set every threshold at $z_j=0.5$ and weight the domains equally at $1/3$.
Profile~A, $(0.55,0.35,0.40)$, clears the threshold on biomedical capacity and
falls below it on the other two, so $d(x)=(0,1,1)$ and $c(x)=2/3$. With a
cut-off of $k=1/2$ it is classified as constrained; with $k=1$ it is not,
because it is not deprived in every domain. Profile~B with structure
unfavourable, $(0.30,0.70,0.20)$, is deprived in biomedical capacity and
structural freedom, so it also scores $2/3$, and the count does not distinguish
the two profiles.

\begin{wecode}[Python]
from hpv.crf import deprivation_count

for name, x in (("Profile A", (0.55, 0.35, 0.40)),
                ("Profile B", (0.30, 0.70, 0.20)),
                ("well placed", (0.90, 0.90, 0.90))):
    c = deprivation_count(x, z=(0.5, 0.5, 0.5))
    print(f"{name:<14} c(x) = {c:.4f}   constrained at k=1/2: {c >= 0.5}")
\end{wecode}

\weanswer{$c(x)=2/3$ for Profile~A and for Profile~B with structure
unfavourable, and $0$ for the well-placed profile. Both reference profiles are
constrained at $k=1/2$ and neither
is at $k=1$. Asserted by \texttt{test\_definition\_deprivation\_count}.}
\end{workedexample}

\begin{reading}[What counting buys and what it costs]
This rule never trades one domain against another. The adequate biomedical
coordinate of Profile~A does not offset its social and structural shortfalls;
it is not counted against them. The cost is that the rule discards magnitude. A
profile at $0.49$ and a profile at $0.05$ are recorded identically, which is unacceptable if the
question is how much help each needs rather than whether either qualifies.

This is the same trade-off as any diagnostic threshold. An HbA1c of $48$
mmol/mol makes a diagnosis; it does not distinguish a patient at $48$ from one
at $110$. Thresholds are for eligibility, continuous scores are for
titration, and neither substitutes for the other.
\end{reading}

The three families are now on the table: additive, CES and threshold. All three
assume the domains can be measured in the first place, and that assumption
deserves its own treatment.

\section{Measurement: Indicators Are Not the Construct}
\label{app:a-measurement}

This section asks how a domain score is obtained from questionnaire answers.
In the first of two models, each
indicator is a signal of one underlying quantity, as several assays of one
analyte each track its concentration.

\begin{definition}[Reflective model]
\[
Y_{ijk}=\tau_{jk}+\lambda_{jk}\eta_{ki}+\varepsilon_{ijk},\qquad
\mathbb E(\varepsilon_{ijk}\mid\eta_{ki})=0.
\]
The latent domain is modelled as a common cause of its indicators.
\end{definition}

In the second, the indicators make up the domain, as the components of a risk
score make up the risk.

\begin{definition}[Formative model]
\[
\eta_{ki}=\sum_jw_{jk}X_{ijk}+\zeta_{ik}.
\]
Indicators are modelled as causes or defining components and need not exhibit
internal consistency.
\end{definition}

\begin{remark}[In plain language]
Reflective indicators behave like effects of an underlying condition;
formative indicators behave like ingredients. Housing and income need not
correlate as symptoms of one disease to remain relevant constraints.
\citep{bollen_lennox_1991,jarvis_2003}
\end{remark}

\begin{reading}[Reflective and formative, at the bedside]
Reflective: a patient has hypothyroidism, and the tiredness, cold intolerance,
constipation and dry skin are four effects of the one underlying state. Treat
the thyroid and all four move together. Asking whether they correlate is a
sensible question, and a symptom that did not correlate with the others would
be evidence it does not belong.

Formative: a patient's cardiovascular risk is built out of age, blood pressure,
lipids and smoking. These do not correlate because a common cause is driving
them; they are the ingredients. A smoker with perfect lipids is unremarkable,
and it would be a mistake to drop smoking from the score because it fails to
correlate with the others.

Biomedical Potential is closer to the first kind, and Systemic Burdens is
closer to the second. Running Cronbach's alpha on the housing, income and
transport indicators of Profile~A and rejecting one for lowering it would be the same error
as dropping smoking from a cardiovascular risk score.
\end{reading}

The next theorem shows that a domain estimated in this way has no natural zero
or unit. Each must be fixed by convention, as an assay reported in arbitrary
units needs a declared reference standard before results can be compared.

\begin{theorem}[Latent scale indeterminacy]
\label{thm:latent_scale}
In $Y=\tau+\lambda\eta+\varepsilon$, define
$\eta^*=a+b\eta$ with $b>0$. The same observed model results from
$\lambda^*=\lambda/b$ and $\tau^*=\tau-\lambda a/b$. Thus the latent
location and scale require identifying constraints.
\end{theorem}
The proof rewrites the model in a shifted and stretched unit and shows that
nothing observable changes.
\begin{proof}
Substitute $\eta=(\eta^*-a)/b$ and collect the intercept and loading terms.
\end{proof}

\begin{remark}[In plain language]
The zero and unit of a latent score are chosen for estimation. Standardising it
does not turn it into years or a natural physical quantity.
\end{remark}

The clinical analogue is the international normalised ratio. The INR is a ratio
constructed precisely so that results are comparable between laboratories using
different thromboplastins; without that deliberate calibration step, a
prothrombin time of $18$ seconds means different things in different
laboratories. A standardised factor score has had no equivalent calibration
performed. A score of $-1.2$ standard deviations for Profile~A states its position in
the NHATS distribution and nothing else. It does not state a loss of $1.2$ of
any quantity.

The next proposition gives the pattern of correlations among indicators that
the higher-order model predicts; Chapter~\ref{ch:methods} tests the fit of that
pattern.

\begin{proposition}[Second-order implied covariance]
\label{prop:secondorder_cov}
If
\[
\mathbf Y=\boldsymbol\tau+\Lambda\boldsymbol\eta+\boldsymbol\varepsilon,
\qquad
\boldsymbol\eta=\boldsymbol\gamma H+\boldsymbol\delta,
\]
$\operatorname{Var}(H)=1$, and the three disturbances are mutually
uncorrelated, then
\[
\Sigma_Y=\Lambda(\boldsymbol\gamma\boldsymbol\gamma^\top+
\Psi_\delta)\Lambda^\top+\Theta_\varepsilon.
\]
\end{proposition}
The proof substitutes one equation into the other.
\begin{proof}
Substitute the second equation into the first and add the covariances of the
mutually uncorrelated terms.
\end{proof}

\begin{remark}[In plain language]
This is the correlation pattern implied by a higher-order factor. It does not
prove such a factor exists. With three first-order domains, the higher-order
layer supplies limited independent evidence; fit alone is insufficient.
\end{remark}

\section{Objects That Must Not Be Confused}
\label{app:a-confusion}

This section separates four quantities that are easily mistaken for one
another.

\begin{definition}[Four distinct objects]
The framework distinguishes:
\begin{enumerate}
 \item the vector profile $X_i=(P_i,S_i,R_i)$;
 \item a chosen normative aggregator $H_{\rho,\alpha}(X_i)$;
 \item a latent score $\widehat H_i=\mathbb E(H_i\mid\mathbf Y_i,\widehat\theta)$;
 \item prognostic risk $\Pr(T_i\le t\mid\widehat H_i,\mathbf Z_i)$.
\end{enumerate}
\end{definition}

The next proposition states that they cannot be swapped.

\begin{proposition}[Non-equivalence]
These objects are not numerically interchangeable without calibration.
\end{proposition}
The proof notes that the four take different kinds of value: three numbers, one
number on a fixed scale, one number on an arbitrary scale, and a probability.
\begin{proof}
They respectively have vector, normalized scalar, arbitrarily identified latent
and time-indexed probability codomains. Their mathematical types and invariance
properties differ.
\end{proof}

\begin{remark}[In plain language]
What a person has, how domains are combined, what a measurement model estimates
and what happens later are different questions. A hazard ratio is not the
definition of H-Potential.
\end{remark}

\begin{workedexample}[apa-four-objects]{The four objects, for Profile~A}
\textbf{(1) The profile.} $X=(0.55,0.35,0.40)$. A description of the domain
values and nothing more. It is a vector, and it has no units.

\textbf{(2) The normative aggregate.} $H_{0,0.5}=0.55\sqrt{0.35\times0.40}
=0.206$. This is the value the chosen ethical rule assigns to the profile.
Change $\rho$ and the number changes while the profile does not: $0.206$ at
$\rho=0$, $0.205$ at $\rho=-1$, $0.193$ in the bottleneck limit. The additive
rule on the same three numbers gives $0.433$, twice as large, because it is a
different rule and not a better measurement.

\textbf{(3) The latent estimate.} If the measurement model returns
$\widehat H=-1.2$ with a standard error of $0.4$, that is what it infers from
the questionnaire answers, on a scale whose zero and unit were chosen for
estimation. Theorem~\ref{thm:latent_scale} says the scale is not identified by
the data, so this number cannot be compared with $0.206$ at all.

\textbf{(4) The predicted risk.} If the hazard model of
Chapter~\ref{ch:results2} returns a five-year mortality probability near
$0.27$, that is a statement about outcomes among people with similar profiles, on
a probability scale, indexed by a time horizon.

\begin{wecode}[Python]
from hpv.crf import H, H_additive

x = (0.55, 0.35, 0.40)
for rho in (0.0, -1.0, -800.0):
    print(f"rho {rho:>7.1f}   H = {H(*x, rho):.4f}")
print("additive rule on the same profile:", round(H_additive(*x), 4))
\end{wecode}

\weanswer{$0.2058$, $0.2053$ and $0.1925$ as $\rho$ falls, against $0.4333$ for
the additive rule. Four numbers, one unchanged profile. Asserted by
\texttt{test\_worked\_example\_four\_distinct\_objects\_for\_profile\_a}.}
\end{workedexample}

\begin{reading}[Which pairs get confused]
Conflating (2) with (3) is the commonest error in composite-index work: a rule
someone chose gets reported as a measurement someone made. Conflating (3) with
(4) is the commonest error in prognostic work: a risk score gets reported as a
description of the person. The framework keeps all four separate for exactly
this reason.
\end{reading}

\section{Failure Conditions and Scope of the Contribution}
\label{app:a-failure}

The next definition lists the findings that would show the single score unfit
for use, as a validation protocol states its rejection criteria before the run.

\begin{definition}[Failure conditions]
The scalar proposal fails in its present form if domain content cannot be
justified independently of outcomes; scales cannot be harmonised; measurement
is unidentified or non-replicable; rankings are unstable over plausible
weights; invariance is inadequate for intended comparisons; or the scalar adds
no interpretable information beyond domain profiles and established measures.
\end{definition}

\begin{remark}[In plain language]
The mathematics allows the theory to fail. If a total score is indefensible,
the profile remains a scientifically legitimate result.
\end{remark}

That is a real possibility and not a rhetorical concession. The most likely
outcome of the empirical chapters is that the domains are measurable, the
higher-order factor is weakly identified, and the profile survives while the
scalar does not. A framework that could not report that outcome would not be
worth building.

Whether the scalar survives or not, the framework eventually has to meet the
question every health measure meets: what does it imply about how money should
be spent? The final section states carefully how little it implies.

\section{Economic Evaluation and Policy Mathematics}
\label{app:a-economics}

This section links the score to decisions about spending. The definition below
gives the standard health-economic measure of value: health gained, less the
health the same money would have bought elsewhere.

\begin{definition}[Net health benefit]
For intervention $a$, willingness-to-pay threshold $\lambda>0$, incremental
health consequence $\Delta Q(a)$ and incremental cost $\Delta C(a)$, define
\[
\operatorname{NHB}(a)=\Delta Q(a)-\frac{\Delta C(a)}{\lambda}.
\]
When $\Delta Q$ is measured in QALYs \citep{whitehead_2010}, NHB is measured in QALYs. The net-benefit form is due to \citet{stinnett_mullahy_1998}.
\end{definition}

\begin{remark}[In plain language]
NHB asks whether the health produced is greater than the health likely displaced
by spending the same resources elsewhere. H-Potential has no threshold
$\lambda$ and cannot replace this calculation.
\end{remark}

\begin{workedexample}[apa-nhb]{Net health benefit for a single intervention}
Suppose an intervention for Profile~B gives $\Delta Q=1.8$ QALYs over the
remaining lifetime and costs $\Delta C=\pounds 540{,}000$ more than the
comparator. At a threshold of $\lambda=\pounds 30{,}000$ per QALY, the health
displaced elsewhere is $540{,}000/30{,}000=18$ QALYs, so
$\operatorname{NHB}=1.8-18=-16.2$ QALYs.

The threshold is the only thing that can change the sign. At
$\lambda=\pounds 300{,}000$ the displaced health is $1.8$ QALYs and
NHB is zero; above that it is positive. Nothing clinical differs
between the two answers.

\begin{wecode}[Python]
from hpv.economics import net_health_benefit, displaced_health

dQ, dC = 1.8, 540_000.0
for lam in (30_000.0, 100_000.0, 300_000.0, 400_000.0):
    print(f"lambda {lam:>9,.0f}  displaced {displaced_health(dC, lam):6.2f} QALYs"
          f"   NHB {net_health_benefit(dQ, dC, lam):+7.2f} QALYs")
\end{wecode}

\weanswer{$\operatorname{NHB}=-16.2$ QALYs at
$\lambda=\pounds 30{,}000$, $-3.6$ at $\pounds 100{,}000$, $0.0$ at
$\pounds 300{,}000$ and $+0.45$ at $\pounds 400{,}000$. Asserted by
\texttt{test\_worked\_example\_treatment\_change\_for\_profile\_b} and
\texttt{test\_definition\_net\_health\_benefit}.}
\end{workedexample}

\begin{reading}[Reading the sign]
The negative sign does not mean the intervention is
ineffective. It produces $1.8$ QALYs. It means that funding it from a fixed
budget removes an estimated $18$ QALYs from people the decision-maker will
never meet. Those people are invisible in the trial and are the whole content
of $\lambda$.

Any change in the H-Potential profile of Profile~B under the intervention
does not enter this calculation. A measure of need and a measure of
opportunity cost answer different questions, and the next proposition states
formally that neither implies the other.
\end{reading}

\begin{proposition}[Need does not imply cost-effectiveness]
A low baseline H-Potential score does not imply
$\operatorname{NHB}(a)>0$ for any particular intervention $a$.
\end{proposition}
The proof needs one example: an intervention that costs money and produces no
health.
\begin{proof}
For any baseline score, choose an intervention with $\Delta Q(a)=0$ and
$\Delta C(a)>0$. Then
$\operatorname{NHB}(a)=-\Delta C(a)/\lambda<0$.
\end{proof}

\begin{remark}[In plain language]
A score can identify serious need without proving that a proposed programme is
effective or a good use of scarce funds. Need, benefit, cost-effectiveness and
affordability are different policy questions.
\end{remark}

The proof is short, but the point matters. Health systems routinely
justify spending by pointing at need, and need is the one thing the
calculation does not use. A low score for Profile~A is a reason to look for an
intervention that helps; it is not evidence that any particular one does. That said,
a rule that maximises net health benefit alone will systematically underserve
people whose disadvantage makes interventions expensive to deliver, which is
why the calculation is extended next rather than left as it stands.

The extension weights health gains by how badly off the group receiving them
is, and one parameter, $\varepsilon$, sets how much.

\begin{definition}[Distribution-sensitive evaluation]
For groups $g=1,\ldots,G$ with population weights $\pi_g$, post-policy health
$q_g(a)>0$ and inequality-aversion \citep{atkinson_1970} parameter $\varepsilon\ge0$, define
\[
W_\varepsilon(a)=
\begin{cases}
\left[\sum_g\pi_gq_g(a)^{1-\varepsilon}\right]^{1/(1-\varepsilon)},
&\varepsilon\ne1,\\
\exp\left(\sum_g\pi_g\log q_g(a)\right),&\varepsilon=1.
\end{cases}
\]
\end{definition}

\begin{remark}[In plain language]
This equally distributed equivalent gives more importance to gains among groups
with worse health as $\varepsilon$ increases. The parameter is a social value
judgement, not an estimated biological fact. Distributional cost-effectiveness
analysis must also include the distribution of opportunity costs.
\citep{asaria_2016_dcea}
\end{remark}

\begin{workedexample}[apa-atkinson]{A transfer between two groups, at three values of $\varepsilon$}
Take two groups of equal size, one with post-policy health $q_1=0.4$ and one
with $q_2=0.8$. At $\varepsilon=0$ the measure is the plain average, $0.600$: a
QALY is a QALY whoever gets it. At $\varepsilon=1$ it is the geometric mean,
$\sqrt{0.4\times0.8}=0.566$. At $\varepsilon=2$ it is the harmonic mean,
$0.533$.

Now move $0.05$ of health from the better-off group to the worse-off, giving
$q_1=0.45$ and $q_2=0.75$. The plain average is unchanged at $0.600$, so at
$\varepsilon=0$ the transfer is worth nothing. At $\varepsilon=1$ the measure
rises to $0.581$ and at $\varepsilon=2$ to $0.5625$. The gap between $0.600$ and
$0.533$ is the price the social judgement is willing to pay for equality, stated
in health rather than hidden in rhetoric.

\begin{wecode}[Python]
from hpv.crf import equity_ede

before, after = [0.4, 0.8], [0.45, 0.75]
for eps in (0.0, 1.0, 2.0):
    b, a = equity_ede(before, eps), equity_ede(after, eps)
    print(f"eps {eps:.0f}   before {b:.4f}   after {a:.4f}   gain {a - b:+.4f}")
\end{wecode}

\weanswer{$0.6000\to0.6000$ at $\varepsilon=0$, $0.5657\to0.5809$ at
$\varepsilon=1$, and $0.5333\to0.5625$ at $\varepsilon=2$. The same transfer is
worth nothing, a little, or a lot, depending only on a declared value judgement.
Asserted by
\texttt{test\_worked\_example\_transfer\_only\_registers\_above\_eps\_zero} and
\texttt{test\_definition\_equity\_ede\_and\_limits}.}
\end{workedexample}

\begin{reading}[Why $\varepsilon$ is declared, not estimated]
The parameter must be declared and varied rather than estimated.
There is no experiment that reveals how much aggregate health a society should
give up to reduce inequality in it.
\end{reading}

The definition below lists the separate pieces of evidence that a funding
decision needs in one health system. They are kept as a list and not combined.

\begin{definition}[Jurisdiction-indexed policy evidence vector]
For policy $a$ in jurisdiction $j$, the minimum decision object is
\[
\mathcal D_j(a)=
(\Delta X_j(a),\Delta Q_j(a),\Delta C_j(a),\operatorname{BI}_j(a),
\operatorname{Dist}_j(a),\operatorname{Unc}_j(a),\operatorname{Gov}_j(a)),
\]
representing profile change, health consequence, cost, budget impact \citep{sullivan_2014},
distribution, uncertainty and governance requirements under the institutions,
prices, population and laws of $j$.
\end{definition}

\begin{remark}[Relation to the eight-coordinate form]
Chapter~\ref{app:crf} carries the same object with the first coordinate
disaggregated, writing $(\Delta P_j,\Delta S_j,\Delta R_j)$ in place of
$\Delta X_j$, which is what makes the decomposition of
\eqref{eq:crf_gap} readable off the vector. That form
omits $\operatorname{Unc}_j$ from the printed list, not because uncertainty is
optional but because it attaches to every other coordinate rather than sitting
beside them: each of $\Delta Q_j$, $\Delta C_j$ and the rest is reported with
its interval. The seven-coordinate form here keeps it explicit; the
eight-coordinate form there distributes it. Neither drops it.
\end{remark}

\begin{remark}[In plain language]
No single coordinate settles the decision. H-Potential contributes
$\Delta X$; it does not absorb safety, effectiveness, costs, affordability,
equity, ethics or accountability into one number.
\end{remark}

The vector is deliberately not collapsed. Collapsing it would produce a single
figure that a committee could rank, and would bury the disagreements that
committee exists to have. A decision about Profile~B can differ between health
systems that measure $\Delta X$ the same way, because they weigh the remaining
six coordinates differently. The following proposition shows that this is a
property of the decision problem, not an inconsistency.

\begin{proposition}[A policy decision is not generally transportable]
Suppose two jurisdictions $j$ and $k$ observe the same H-Potential profile
change, $\Delta X_j(a)=\Delta X_k(a)$, and the same health consequence,
$\Delta Q_j(a)=\Delta Q_k(a)$. These equalities are insufficient to imply the
same policy ranking of $a$.
\end{proposition}

The proof builds two health systems with the same effect and opposite
decisions.
\begin{proof}
Let a simplified jurisdictional net-benefit rule be
$N_j(a)=\lambda_j\Delta Q_j(a)-\Delta C_j(a)$, subject also to a budget
constraint. Equality of $\Delta X$ and $\Delta Q$ imposes no equality on the
opportunity-cost parameter $\lambda_j$, incremental cost $\Delta C_j$, or
budget impact $\operatorname{BI}_j$. Choose admissible values for which
$N_j(a)>0$ and $N_k(a)<0$. The rankings differ despite equal profile and health
effects. Distributional weights or legal constraints provide further grounds
for divergence. Hence transport of effect does not entail transport of
decision.
\end{proof}

\begin{remark}[In plain language]
An intervention can work equally well in two places and still be a sound use of
resources in only one of them. What must be given up, who bears that loss, what
the public budget can sustain and which rights the state must honour are local
facts and judgements. The mathematics therefore forbids a global
``cost-per-H-Potential'' threshold.
\end{remark}

The results proved here are conditional properties of stated forms. They do not
prove that the domains exhaust health capability, that CES is a biological law,
that any weights are ethically correct, or that a common latent factor exists.
The contribution is a disciplined bridge: a domain space grounded in
healthy-ageing and capability theory; axioms for any scalar; explicit
compensatory and non-compensatory families; proofs of their implications; and a
measurement architecture able to reject aggregation.


\section*{What Is New in This Appendix}

Before this appendix, the separate pieces existed: measurement scales and the
transformations they permit \citep{stevens_1946}, conditions for representing
an ordering by a number \citep{debreu_1954}, additive multi-attribute value
\citep{keeney_raiffa_1976}, the CES family \citep{arrow_chenery_1961}, counting
of deprivations (\S\ref{app:a-threshold}), reflective and formative indicators
\citep{bollen_lennox_1991,jarvis_2003}, net health benefit
\citep{stinnett_mullahy_1998} and inequality-averse summaries
\citep{atkinson_1970}. None was written for a health-capability score.
This appendix adds four things.
It states the requirements for such a score before choosing its form, and names
the one a reader may reject, the capacity ceiling $H\le p$. It sets a relational
score beside the ceiling score and proves that fit to data cannot choose between
them. It separates the profile, the chosen score, the estimated score and the
predicted risk, and proves that they are not interchangeable. It proves that
need does not imply value for money and that a funding decision does not carry
over between health systems. Each claim can be checked in two ways: each proof
is a few lines of algebra or a single counterexample, and each worked number is
asserted by a named test in the verification suite whose run
Appendix~\ref{app:verification} reports.
